\documentclass[12pt,a4paper]{article}

% ─── Packages ────────────────────────────────────────────────────────────────
\usepackage[utf8]{inputenc}
\usepackage[T1]{fontenc}
\usepackage[french]{babel}
\usepackage{geometry}
\geometry{margin=2.5cm}

\usepackage{amsmath, amssymb, amsthm}
\usepackage{algorithm}
\usepackage{algpseudocode}
\usepackage{tikz}
\usetikzlibrary{arrows.meta, positioning, shapes.geometric, decorations.markings}
\usepackage{pgfplots}
\pgfplotsset{compat=1.18}
\usepackage{booktabs}
\usepackage{array}
\usepackage{xcolor}
\usepackage{hyperref}
\usepackage{enumitem}
\usepackage{mdframed}
\usepackage{lmodern}
\usepackage{microtype}
\usepackage{graphicx}
\usepackage{caption}
\usepackage{subcaption}
\usepackage{tcolorbox}
\tcbuselibrary{skins, breakable}

% ─── Colors ──────────────────────────────────────────────────────────────────
\definecolor{urgence}{RGB}{180,30,30}
\definecolor{theorie}{RGB}{30,80,160}
\definecolor{pratique}{RGB}{20,130,80}
\definecolor{attention}{RGB}{200,120,0}
\definecolor{grisléger}{RGB}{245,245,245}
\definecolor{bleuléger}{RGB}{230,240,255}
\definecolor{vertléger}{RGB}{230,248,238}

% ─── Custom environments ──────────────────────────────────────────────────────
\newtcolorbox{defbox}[1]{
  colback=bleuléger, colframe=theorie, fonttitle=\bfseries,
  title={#1}, breakable, sharp corners=south
}

\newtcolorbox{metierbox}[1]{
  colback=vertléger, colframe=pratique, fonttitle=\bfseries,
  title={#1}, breakable, sharp corners=south
}

\newtcolorbox{alertbox}[1]{
  colback=yellow!10, colframe=attention, fonttitle=\bfseries,
  title={#1}, breakable, sharp corners=south
}

\newtcolorbox{jurybox}[1]{
  colback=red!5, colframe=urgence, fonttitle=\bfseries,
  title={#1}, breakable, sharp corners=south
}

% ─── Theorem styles ───────────────────────────────────────────────────────────
\theoremstyle{definition}
\newtheorem{definition}{Définition}[section]
\newtheorem{propriete}{Propriété}[section]
\newtheorem{invariant}{Invariant métier}[section]

\theoremstyle{remark}
\newtheorem{remarque}{Remarque}[section]
\newtheorem{exemple}{Exemple}[section]

% ─── Hyperref setup ───────────────────────────────────────────────────────────
\hypersetup{
  colorlinks=true,
  linkcolor=theorie,
  urlcolor=urgence,
  pdftitle={Routage d'urgence dynamique avec A*},
  pdfauthor={Système de routage critique},
}

% ─────────────────────────────────────────────────────────────────────────────
\begin{document}

% ═══════════════════════════════════════════════════════════════════
%  COUVERTURE
% ═══════════════════════════════════════════════════════════════════
\begin{titlepage}
  \centering
  \vspace*{2cm}
  {\color{urgence}\rule{\linewidth}{2pt}}\\[0.5em]
  {\Huge\bfseries Moteur de Routage d'Urgence}\\[0.4em]
  {\Large\bfseries Algorithme A* à Poids Dynamiques}\\[0.4em]
  {\color{urgence}\rule{\linewidth}{2pt}}\\[2em]

  \begin{tikzpicture}
    \node[circle, fill=urgence!20, draw=urgence, thick, minimum size=2cm] (s) at (0,0) {\textbf{S}};
    \node[circle, fill=theorie!20, draw=theorie, thick, minimum size=1.5cm] (a) at (3,1.5) {A};
    \node[circle, fill=theorie!20, draw=theorie, thick, minimum size=1.5cm] (b) at (3,-1.5) {B};
    \node[circle, fill=theorie!20, draw=theorie, thick, minimum size=1.5cm] (c) at (6,0) {C};
    \node[circle, fill=pratique!20, draw=pratique, thick, minimum size=2cm] (t) at (9,0) {\textbf{T}};
    \draw[-{Latex[length=3mm]}, theorie, thick] (s) -- node[above,sloped,font=\small]{$w=3$} (a);
    \draw[-{Latex[length=3mm]}, urgence, thick, dashed] (s) -- node[below,sloped,font=\small]{$w=5\to8$} (b);
    \draw[-{Latex[length=3mm]}, theorie, thick] (a) -- node[above,sloped,font=\small]{$w=2$} (c);
    \draw[-{Latex[length=3mm]}, theorie, thick] (b) -- node[below,sloped,font=\small]{$w=4$} (c);
    \draw[-{Latex[length=3mm]}, pratique, thick] (c) -- node[above,sloped,font=\small]{$w=1$} (t);
    \node[font=\footnotesize, color=urgence] at (1.5,-2.5) {Arc vieilli (trafic)};
  \end{tikzpicture}\\[2em]

  {\large Note technique — Systèmes critiques d'aide à la décision}\\[1em]
  {\large\color{gray} Contexte : services d'urgence (SAMU, pompiers, ambulances)}\\[3em]

  \begin{tabular}{rl}
    \textbf{Domaine :} & Algorithmique appliquée aux systèmes critiques \\
    \textbf{Données :} & Graphe OSM Yaoundé + trafic horaire par segment \\
    \textbf{Révision :} & Poids actualisés toutes les 5 minutes \\
    \textbf{Véhicules :} & Ambulance, camion de pompiers, véhicule léger \\
  \end{tabular}

  \vfill
  {\small\color{gray} Document pédagogique — Présentation jury et équipe technique}
\end{titlepage}

% ═══════════════════════════════════════════════════════════════════
%  TABLE DES MATIÈRES
% ═══════════════════════════════════════════════════════════════════
\tableofcontents
\newpage

% ═══════════════════════════════════════════════════════════════════
%  INTRODUCTION
% ═══════════════════════════════════════════════════════════════════
\section*{Introduction}
\addcontentsline{toc}{section}{Introduction}

Un moteur de routage pour les services d'urgence doit prendre des décisions
\emph{rapides}, \emph{fiables} et \emph{adaptatives}. Contrairement à la
navigation civile classique (Google Maps, Waze), le routage d'urgence opère
dans un environnement où :

\begin{itemize}
  \item chaque minute perdue peut coûter une vie ;
  \item les conditions de circulation changent en continu (accidents, pluie,
        barrages, foules) ;
  \item les contraintes de passabilité dépendent du véhicule (ambulance
        $\neq$ camion de pompiers $\neq$ véhicule léger) ;
  \item le système doit être \emph{auditabl}e et \emph{explicable} devant
        un jury ou une autorité de certification.
\end{itemize}

Ce document explique, de manière progressive, comment l'algorithme A*
peut être adapté à ce contexte, comment les poids des arcs évoluent dans le
temps (\emph{vieillissement}), et quels invariants métier le système doit
respecter en toute circonstance.

\bigskip
\noindent\textbf{Données réelles utilisées.}
Les exemples numériques s'appuient sur deux jeux de données concrets :
\begin{enumerate}
  \item \texttt{export.geojson} — graphe routier OpenStreetMap (Yaoundé)
        contenant arcs, nœuds et attributs de voiries ;
  \item \texttt{traffic\_by\_segment\_hour.csv} — 2\,400 enregistrements
        de trafic horaire par segment, incluant vitesse moyenne, niveau de
        congestion, multiplicateur de temps de parcours, passabilité ambulance
        / camion, et présence de pluie.
\end{enumerate}

\newpage

% ═══════════════════════════════════════════════════════════════════
%  SECTION 1 — THÉORIE DES GRAPHES ET POIDS DYNAMIQUES
% ═══════════════════════════════════════════════════════════════════
\section{Graphes pondérés et vieillissement des arcs}

\subsection{Modélisation du réseau routier comme graphe pondéré}

\begin{defbox}{Définition : Graphe pondéré orienté}
Un \textbf{graphe pondéré orienté} est un triplet $G = (V, E, w)$ où :
\begin{itemize}
  \item $V$ est l'ensemble des \textbf{nœuds} (intersections, points remarquables) ;
  \item $E \subseteq V \times V$ est l'ensemble des \textbf{arcs} (segments de route) ;
  \item $w : E \to \mathbb{R}^+$ est une fonction de \textbf{pondération} qui
        associe à chaque arc $(u,v)$ un coût $w(u,v) > 0$.
\end{itemize}
\end{defbox}

Dans le contexte routier :
\begin{center}
\begin{tabular}{lll}
  \toprule
  \textbf{Élément du graphe} & \textbf{Objet réel} & \textbf{Exemple} \\
  \midrule
  Nœud $v \in V$ & Intersection, rond-point & Carrefour Nlongkak \\
  Arc $(u,v) \in E$ & Segment de voie & Tronçon avenue Kennedy \\
  Poids $w(u,v)$ & Temps de parcours (secondes) & 47\,s en heure creuse \\
  \bottomrule
\end{tabular}
\end{center}

\begin{remarque}
  Le poids $w(u,v)$ n'est \emph{pas} une distance physique, mais un
  \textbf{temps de parcours estimé}, dépendant de la vitesse, du trafic,
  des conditions météo et du type de véhicule.
\end{remarque}

\subsection{Le vieillissement des arcs — définition formelle}

\begin{defbox}{Définition : Arc temporellement pondéré (poids vieillissant)}
On dit qu'un arc $(u,v)$ est \textbf{vieillissant} si son poids est une
fonction du temps :
$$
w(u,v,t) = w_0(u,v) \cdot \mu(u,v,t)
$$
où :
\begin{itemize}
  \item $w_0(u,v)$ est le \textbf{poids de base} (temps à vitesse libre,
        conditions idéales) ;
  \item $\mu(u,v,t) \geq 1$ est le \textbf{multiplicateur temporel},
        reflétant la dégradation des conditions à l'instant $t$.
\end{itemize}
\end{defbox}

\noindent\textbf{Données réelles — multiplicateurs observés} (fichier
\texttt{traffic\_by\_segment\_hour.csv}, 2\,400 relevés) :

\begin{center}
\begin{tabular}{lrrr}
  \toprule
  \textbf{Type de voie} & \textbf{Multiplicateur moyen} &
  \textbf{Max observé} & \textbf{Interprétation} \\
  \midrule
  Autoroute / trunk  & $\times 1{,}57$  & $\times 3{,}33$ &
      Congestion possible aux heures de pointe \\
  Primaire           & $\times 1{,}58$  & $\times 3{,}33$ &
      Axes urbains soumis à la congestion \\
  Secondaire         & $\times 1{,}77$  & $\times 3{,}33$ &
      Voies mixtes, trafic irrégulier \\
  Tertiaire          & $\times 1{,}78$  & $\times 3{,}33$ &
      Routes de quartier, congestion fréquente \\
  Résidentielle      & $\times 1{,}78$  & $\times 3{,}33$ &
      Voies étroites, imprévus fréquents \\
  \bottomrule
\end{tabular}
\end{center}

\begin{alertbox}{Interprétation critique}
Un multiplicateur maximal de $\times 3{,}33$ signifie qu'un trajet de
10 minutes en conditions idéales peut prendre \textbf{33 minutes} en
situation de congestion extrême. Pour une ambulance, cette différence
est potentiellement fatale.
\end{alertbox}

\subsection{Illustration graphique du vieillissement}

\begin{figure}[h]
\centering
\begin{tikzpicture}
  \begin{axis}[
    width=13cm, height=6cm,
    xlabel={Heure de la journée},
    ylabel={Multiplicateur $\mu(t)$},
    xmin=0, xmax=23,
    ymin=1, ymax=3.5,
    xtick={0,3,6,9,12,15,18,21},
    grid=major,
    grid style={gray!30},
    legend pos=north west,
    title={Évolution temporelle du poids d'un arc — exemple autoroute vs résidentiel}
  ]
  % Motorway approximation
  \addplot[color=theorie, thick, smooth] coordinates {
    (0,1.3) (3,1.2) (6,1.8) (7,2.5) (8,2.8) (9,2.2)
    (12,1.6) (13,1.5) (17,2.4) (18,3.0) (19,2.5) (21,1.5) (23,1.2)
  };
  \addlegendentry{Autoroute ($\mu$ moyen $1{,}57$)}

  % Residential approximation
  \addplot[color=urgence, thick, smooth, dashed] coordinates {
    (0,1.4) (3,1.3) (6,2.0) (7,2.8) (8,3.2) (9,2.6)
    (12,1.9) (13,1.7) (17,2.7) (18,3.3) (19,2.9) (21,1.8) (23,1.4)
  };
  \addlegendentry{Résidentielle ($\mu$ moyen $1{,}78$)}

  % 5-minute window
  \addplot[fill=attention!20, draw=none, forget plot]
    coordinates {(8,1)} -- (axis cs:8,3.5) -- (axis cs:8.083,3.5) -- (axis cs:8.083,1) -- cycle;
  \node[font=\tiny, color=attention] at (axis cs:8.5, 3.3) {fenêtre 5 min};
  \end{axis}
\end{tikzpicture}
\caption{Le poids d'un arc varie au cours du temps. La fenêtre de 5 minutes
  représente l'intervalle de recalcul du chemin optimal.}
\end{figure}

\subsection{Pourquoi un chemin optimal à $t_0$ devient sous-optimal à $t_0 + \Delta t$}

\begin{defbox}{Propriété : Instabilité temporelle de l'optimalité}
Soit $P^*(t_0) = (v_0, v_1, \ldots, v_k)$ le chemin optimal calculé à $t_0$.
Son coût total est :
$$
\text{coût}(P^*, t_0) = \sum_{i=0}^{k-1} w(v_i, v_{i+1}, t_0)
$$
À $t_0 + \Delta t$, ce même chemin a un coût réel :
$$
\text{coût}(P^*, t_0 + \Delta t) = \sum_{i=0}^{k-1} w(v_i, v_{i+1},\; t_0 + \Delta t)
$$
Il existe un chemin alternatif $P'$ tel que :
$$
\text{coût}(P', t_0 + \Delta t) < \text{coût}(P^*, t_0 + \Delta t)
$$
$P^*$ n'est alors \textbf{plus optimal}.
\end{defbox}

\paragraph{Scénario concret.}
Une ambulance part à $t_0 =$ 08h00 par l'avenue Kennedy (chemin optimal).
À 08h05, un accident bloque cette avenue — le multiplicateur passe de
$\times 1{,}8$ à $\times 3{,}2$. Sans recalcul, l'ambulance s'engage
dans un chemin devenu 78\,\% plus lent que l'alternative disponible.

\bigskip
\noindent C'est précisément pourquoi le moteur recalcule le chemin \textbf{toutes
les 5 minutes} (ou à chaque événement critique détecté).

\newpage

% ═══════════════════════════════════════════════════════════════════
%  SECTION 2 — L'ALGORITHME A* ET SON ADAPTATION DYNAMIQUE
% ═══════════════════════════════════════════════════════════════════
\section{L'algorithme A* et son adaptation aux poids dynamiques}

\subsection{Rappel : A* classique}

A* est un algorithme de recherche de chemin le plus court qui améliore
Dijkstra en utilisant une \textbf{heuristique} pour guider l'exploration
vers la destination, réduisant ainsi le nombre de nœuds explorés.

\begin{defbox}{Fonction d'évaluation de A*}
Pour chaque nœud $n$ exploré, A* calcule :
$$
f(n) = g(n) + h(n)
$$
\begin{itemize}
  \item $g(n)$ : coût réel du chemin depuis la source jusqu'à $n$ ;
  \item $h(n)$ : estimation heuristique du coût restant de $n$ jusqu'à
        la destination ;
  \item $f(n)$ : estimation du coût total du chemin passant par $n$.
\end{itemize}
A* explore en priorité les nœuds avec la plus petite valeur de $f(n)$.
\end{defbox}

\begin{propriete}[Admissibilité de l'heuristique]
  Pour qu'A* garantisse l'optimalité, l'heuristique $h(n)$ doit être
  \textbf{admissible} : elle ne doit jamais surestimer le coût réel.
  $$
  h(n) \leq h^*(n) \quad \forall n \in V
  $$
  où $h^*(n)$ est le coût réel optimal depuis $n$ vers la destination.
\end{propriete}

\noindent\textbf{Heuristique recommandée pour le routage urbain :}
$$
h(n) = \frac{d_{\text{euclidienne}}(n, \text{destination})}{v_{\max}}
$$
où $v_{\max}$ est la vitesse maximale autorisée sur le réseau.
Cette heuristique est admissible car un véhicule ne peut pas aller plus
vite que $v_{\max}$.

\subsection{Pseudocode d'A* adapté aux poids dynamiques}

\begin{algorithm}[H]
\caption{A* dynamique avec poids temporels et contraintes véhicule}
\begin{algorithmic}[1]
\Require Graphe $G=(V,E,w)$, source $s$, destination $t$,
         instant $t_0$, type véhicule $\tau$
\Ensure Chemin optimal $P^*$ ou \textsc{Aucun}

\State $\text{ouvert} \gets \{s\}$ \Comment{File de priorité}
\State $g[s] \gets 0$
\State $f[s] \gets h(s, t)$
\State $\text{parent}[s] \gets \text{null}$

\While{$\text{ouvert} \neq \emptyset$}
  \State $n \gets \arg\min_{x \in \text{ouvert}} f[x]$
  \If{$n = t$}
    \Return \Call{ReconstruireChemin}{parent, $t$}
  \EndIf
  \State Retirer $n$ de ouvert

  \For{chaque voisin $m$ de $n$, arc $(n,m)$}
    \If{\textbf{non} \Call{Passable}{$(n,m)$, $\tau$, $t_0$}}
      \State \textbf{continuer} \Comment{Filtre de passabilité véhicule}
    \EndIf
    \State $w_{\text{eff}} \gets w_0(n,m) \cdot \mu(n,m, t_0) \cdot \phi_\tau(n,m)$
    \State $g_{\text{tentative}} \gets g[n] + w_{\text{eff}}$
    \If{$g_{\text{tentative}} < g[m]$}
      \State $g[m] \gets g_{\text{tentative}}$
      \State $f[m] \gets g[m] + h(m, t)$
      \State $\text{parent}[m] \gets n$
      \State Ajouter $m$ à ouvert (ou mettre à jour sa priorité)
    \EndIf
  \EndFor
\EndWhile
\Return \textsc{Aucun} \Comment{Destination inatteignable}
\end{algorithmic}
\end{algorithm}

\paragraph{Nouveaux paramètres par rapport à A* classique :}
\begin{itemize}
  \item $\mu(n,m,t_0)$ : multiplicateur de congestion à l'instant $t_0$
        (issu du fichier de trafic) ;
  \item $\phi_\tau(n,m)$ : facteur de pénalité propre au type de véhicule $\tau$
        (ambulance, camion de pompiers, etc.) ;
  \item \textsc{Passable}$(e, \tau, t)$ : prédicat booléen vérifiant si
        l'arc $e$ est physiquement franchissable par le véhicule $\tau$
        à l'instant $t$.
\end{itemize}

\subsection{Modélisation des facteurs de pondération}

\subsubsection{Facteur de congestion $\mu(u,v,t)$}

Le multiplicateur $\mu$ est extrait de la table de trafic horaire et
interpolé linéairement entre deux mesures consécutives :

$$
\mu(u,v,t) = \mu_h + \frac{t - t_h}{t_{h+1} - t_h} \cdot (\mu_{h+1} - \mu_h)
$$

où $h = \lfloor t / 3600 \rfloor$ est l'heure courante.

\subsubsection{Facteur météorologique $\rho(u,v,t)$}

La pluie est codifiée dans les données (\texttt{rain = oui/non}).
En présence de pluie, le poids est pénalisé :

$$
\rho(u,v,t) = \begin{cases} 1{,}0 & \text{si sec} \\ 1{,}25 & \text{si pluie légère} \\ 1{,}6 & \text{si pluie forte} \end{cases}
$$

\subsubsection{Facteur de passabilité par type de véhicule $\phi_\tau$}

\begin{center}
\begin{tabular}{lccc}
  \toprule
  \textbf{Type de voie} & \textbf{Ambulance} &
  \textbf{Camion pompiers} & \textbf{Véh. léger} \\
  \midrule
  Autoroute    & $\phi = 1{,}0$ & $\phi = 1{,}0$ & $\phi = 1{,}0$ \\
  Primaire     & $\phi = 1{,}0$ & $\phi = 1{,}05$ & $\phi = 1{,}0$ \\
  Secondaire   & $\phi = 1{,}05$ & $\phi = 1{,}15$ & $\phi = 1{,}0$ \\
  Tertiaire    & $\phi = 1{,}1$ & $\phi = 1{,}3$ & $\phi = 1{,}0$ \\
  Résidentielle & $\phi = 1{,}15$ & \textit{bloqué} & $\phi = 1{,}0$ \\
  Chemin piéton & \textit{bloqué} & \textit{bloqué} & \textit{bloqué} \\
  \bottomrule
\end{tabular}
\end{center}

\noindent La valeur \emph{bloqué} est implémentée en affectant
$w_{\text{eff}} = +\infty$ (ou en sautant l'arc via le prédicat
\textsc{Passable}).

\subsubsection{Poids effectif complet}

Le poids effectif final utilisé par A* est :

\begin{equation}
\boxed{
  w_{\text{eff}}(u,v,t,\tau) =
    w_0(u,v) \cdot \mu(u,v,t) \cdot \rho(u,v,t) \cdot \phi_\tau(u,v)
}
\end{equation}

\subsection{Recalcul périodique — stratégie de réinitialisation}

Le recalcul toutes les 5 minutes n'implique pas nécessairement de
repartir de zéro. Deux stratégies sont envisageables :

\begin{center}
\begin{tabular}{lll}
  \toprule
  \textbf{Stratégie} & \textbf{Avantage} & \textbf{Inconvénient} \\
  \midrule
  Recalcul complet (A* from scratch) &
      Toujours optimal & Coûteux si graphe large \\
  Recalcul incrémental (D* Lite) &
      Rapide, ne retraite que les arcs modifiés &
      Implémentation complexe \\
  Recalcul partiel (rayon $r$ km) &
      Bon compromis & Peut rater une déviation lointaine \\
  \bottomrule
\end{tabular}
\end{center}

\begin{metierbox}{Recommandation pour les services d'urgence}
  Pour les systèmes temps-réel avec des graphes de taille urbaine
  ($|V| < 100\,000$ nœuds), le \textbf{recalcul complet toutes les 5
  minutes} est préférable pour sa \emph{simplicité}, sa
  \emph{correctness} garantie, et sa \emph{auditabilité}.
  La latence de recalcul est typiquement inférieure à 200\,ms sur
  un réseau urbain standard.
\end{metierbox}

\newpage

% ═══════════════════════════════════════════════════════════════════
%  SECTION 3 — NÉCESSITÉ DU RECALCUL PÉRIODIQUE
% ═══════════════════════════════════════════════════════════════════
\section{Nécessité du recalcul périodique}

\subsection{Le problème du chemin fossilisé}

\begin{defbox}{Définition : Chemin fossilisé}
  Un \textbf{chemin fossilisé} est un chemin calculé à $t_0$ et suivi
  sans recalcul jusqu'à $t_0 + T$, alors que le graphe a évolué entre
  ces deux instants. Il peut devenir largement sous-optimal ou même
  infranchissable.
\end{defbox}

\paragraph{Exemple chiffré.}
Supposons un trajet source $\to$ destination de distance $d = 8$\,km
sur autoroute. À $t_0$, le multiplicateur est $\mu = 1{,}2$.
À $t_0 + 10$\,min, un accident survient et $\mu$ passe à $3{,}0$.
Une voie alternative secondaire avec $\mu = 1{,}5$ existait.

\begin{center}
\begin{tabular}{lccc}
  \toprule
  & \textbf{Chemin fossilisé} & \textbf{Alternative} &
  \textbf{Écart} \\
  \midrule
  Vitesse libre & 90 km/h & 60 km/h & --- \\
  Multiplicateur réel & $\times 3{,}0$ & $\times 1{,}5$ & --- \\
  Vitesse effective & 30 km/h & 40 km/h & --- \\
  Temps de parcours & 16 min & 12 min & \textbf{+4 min} \\
  \bottomrule
\end{tabular}
\end{center}

\noindent 4 minutes de retard pour une ambulance en arrêt cardiaque peuvent
représenter une différence de survie significative (la survie au
défibrillateur décroît d'environ 10\,\% par minute).

\subsection{Événements déclencheurs d'un recalcul anticipé}

En plus du recalcul périodique toutes les 5 minutes, certains événements
doivent déclencher un recalcul \emph{immédiat} :

\begin{enumerate}
  \item \textbf{Détection d'un arc bloqué} (accident, travaux) sur le
        chemin courant ;
  \item \textbf{Hausse brutale du multiplicateur} d'un arc du chemin
        courant ($\Delta\mu > \theta$, seuil à définir) ;
  \item \textbf{Changement de conditions météo} (début de pluie forte) ;
  \item \textbf{Signal du conducteur} (déviation manuelle, obstacle visuel).
\end{enumerate}

\subsection{Schéma de la boucle de recalcul}

\begin{figure}[h]
\centering
\begin{tikzpicture}[
  box/.style={rectangle, rounded corners, draw=theorie, fill=bleuléger,
               minimum width=3.5cm, minimum height=1cm, align=center, font=\small},
  urgbox/.style={rectangle, rounded corners, draw=urgence, fill=red!8,
               minimum width=3.5cm, minimum height=1cm, align=center, font=\small},
  arr/.style={-{Latex[length=3mm]}, thick}
]
  \node[box] (init)  at (0,0)    {Départ mission\\à $t_0$};
  \node[box] (calc)  at (4,0)    {Calcul A*\\$P^*(t_0)$};
  \node[box] (nav)   at (8,0)    {Navigation\\sur $P^*$};
  \node[box] (timer) at (8,-2.5) {Minuterie\\5 min écoulées ?};
  \node[urgbox] (evt) at (4,-2.5) {Événement\\critique ?};
  \node[box] (update) at (0,-2.5) {Mise à jour\\des poids $w(t)$};
  \node[box] (arrive) at (8,2)   {Arrivée\\destination};

  \draw[arr] (init) -- (calc);
  \draw[arr] (calc) -- (nav);
  \draw[arr] (nav) -- (arrive);
  \draw[arr] (nav) -- (timer);
  \draw[arr] (timer) -- node[right, font=\tiny]{oui} (evt);
  \draw[arr] (timer.east) -- ++(0.8,0) |- node[right, font=\tiny, pos=0.25]{non} (nav.east);
  \draw[arr] (evt) -- node[above, font=\tiny]{oui} (update);
  \draw[arr] (evt.east) -- ++(0.5,0) |- node[right, font=\tiny, pos=0.25]{non} (nav.east);
  \draw[arr] (update) -- (calc);
\end{tikzpicture}
\caption{Boucle de recalcul périodique et événementiel du moteur de routage.}
\end{figure}

\newpage

% ═══════════════════════════════════════════════════════════════════
%  SECTION 4 — INVARIANTS MÉTIER
% ═══════════════════════════════════════════════════════════════════
\section{Invariants métier et contraintes de sécurité}

Un invariant métier est une \textbf{propriété qui doit rester vraie à tout
instant}, indépendamment des optimisations algorithmiques.

\subsection{Invariants fondamentaux}

\begin{invariant}[Accessibilité garantie]
  Le système doit \textbf{toujours} fournir un chemin vers la destination
  si la destination est physiquement accessible, même si ce chemin n'est
  pas optimal.
  $$
  \exists\, P : s \xrightarrow{*} t \;\Rightarrow\; \text{moteur retourne } P \neq \emptyset
  $$
\end{invariant}

\begin{invariant}[Passabilité physique du véhicule]
  Aucun arc retourné dans le chemin ne doit être marqué comme impraticable
  pour le type de véhicule $\tau$ demandeur.
  $$
  \forall (u,v) \in P^* : \textsc{Passable}(u,v,\tau,t) = \text{vrai}
  $$
\end{invariant}

\begin{invariant}[Délai maximal admissible]
  Le chemin retourné doit avoir un coût (temps de parcours estimé)
  inférieur au délai maximal $D_{\max}$ défini par le protocole médical.
  $$
  \text{coût}(P^*) \leq D_{\max}
  $$
  Si aucun chemin ne respecte $D_{\max}$, le système doit émettre une
  \textbf{alerte critique} et proposer le chemin le plus court disponible.
\end{invariant}

\begin{invariant}[Continuité de service]
  Un recalcul en cours ne doit pas laisser le véhicule sans instruction
  de navigation. La route en cours reste active jusqu'à la disponibilité
  du nouveau chemin.
  $$
  \forall t : P_{\text{actif}}(t) \neq \emptyset
  $$
\end{invariant}

\begin{invariant}[Admissibilité de l'heuristique]
  L'heuristique $h$ utilisée par A* doit rester admissible après chaque
  mise à jour des poids, pour garantir l'optimalité du résultat.
  $$
  h(n) \leq \frac{d_{\text{euclidienne}}(n, t)}{v_{\max}} \quad \forall n,t
  $$
\end{invariant}

\subsection{Contraintes de sécurité opérationnelle}

\begin{center}
\begin{tabular}{lll}
  \toprule
  \textbf{Contrainte} & \textbf{Valeur recommandée} & \textbf{Raison} \\
  \midrule
  Délai max. recalcul & $< 200$\,ms & Temps réel conducteur \\
  Fréquence recalcul & toutes les 5\,min & Compromis fraîcheur/charge \\
  Âge max. des données & $< 15$\,min & Données périmées = risque \\
  Poids minimal arc & $> 0$ & Éviter les boucles infinies \\
  Multiplicateur plafond & $\leq 5{,}0$ & Éviter les divergences numériques \\
  Hauteur de flot ambulance & voie $\geq 3{,}5$\,m & Contrainte physique \\
  \bottomrule
\end{tabular}
\end{center}

\subsection{Gestion des données périmées (stale data)}

\begin{alertbox}{Risque : Données de trafic périmées}
Si la source de données trafic est indisponible (panne réseau, timeout),
le système doit basculer sur une \textbf{stratégie de repli} :
\begin{enumerate}
  \item Utiliser les données historiques moyennes (moyenne de la même
        heure sur les 30 derniers jours) ;
  \item Augmenter conservativement tous les multiplicateurs de $+20\,\%$
        pour compenser l'incertitude ;
  \item Signaler visuellement à l'opérateur que les données sont
        estimées, pas mesurées.
\end{enumerate}
\end{alertbox}

\newpage

% ═══════════════════════════════════════════════════════════════════
%  SECTION 5 — PRÉSENTATION AU JURY NON TECHNIQUE
% ═══════════════════════════════════════════════════════════════════
\section{Communication au jury et aux équipes non techniques}

\subsection{La métaphore du GPS qui écoute la radio}

\begin{jurybox}{Message clé pour le jury}
  \textbf{Notre moteur de routage fonctionne comme un GPS
  qui écoute la radio en permanence.}

  \medskip
  Un GPS classique calcule votre itinéraire une seule fois
  et vous guide en ligne droite, même si un accident vient de bloquer
  la route. Notre système, lui, \textbf{réécoute la situation toutes
  les 5 minutes} et se demande : \og{}Est-ce que le chemin actuel
  est toujours le plus rapide ?\fg{} Si non, il recalcule et propose
  une déviation, automatiquement, sans intervention humaine.
\end{jurybox}

\subsection{Les trois questions que le jury posera}

\paragraph{Question 1 : \og{}Et si le système se trompe ?\fg{}}

\noindent\textbf{Réponse :} Le système ne se trompe pas sur les données
disponibles — il calcule l'optimal compte tenu de ce qu'il sait. Si les
données sont imparfaites (capteur défaillant, incident non signalé), nous
avons un mécanisme de repli conservatif qui surestime délibérément les
temps de trajet pour rester en sécurité. De plus, le conducteur garde
\emph{toujours} la main : l'IA suggère, l'humain décide.

\paragraph{Question 2 : \og{}Pourquoi ne pas garder le même chemin jusqu'au bout ?\fg{}}

\noindent\textbf{Réponse :} Parce que la ville est vivante. En 10 minutes,
un feu de circulation tombe en panne, un camion se gare en double file,
la pluie rend une route glissante. Un chemin calculé à 08h00 peut devenir
33\,\% plus lent à 08h10. Pour une ambulance, ces minutes comptent.

\paragraph{Question 3 : \og{}Comment savez-vous que le nouveau chemin est vraiment meilleur ?\fg{}}

\noindent\textbf{Réponse :} L'algorithme A* est prouvé mathématiquement
optimal : si l'heuristique est admissible (ce que nous garantissons), il
trouve \emph{toujours} le chemin le moins coûteux parmi tous les chemins
possibles. Ce n'est pas une heuristique floue — c'est une garantie formelle.

\subsection{Tableau de bord de transparence}

Pour assurer la transparence du système, chaque décision de routage
doit être accompagnée d'un log structuré :

\begin{center}
\begin{tabular}{ll}
  \toprule
  \textbf{Champ du log} & \textbf{Exemple de valeur} \\
  \midrule
  Timestamp de calcul & 2024-11-15 08:07:32 \\
  Véhicule & Ambulance SAMU-12 \\
  Chemin retourné & [N1043 $\to$ A22 $\to$ R47 $\to$ N8] \\
  Coût estimé & 14 min 20 s \\
  Multiplicateurs appliqués & $\mu_1 = 1.8,\; \mu_2 = 1.2,\; \mu_3 = 2.1$ \\
  Pluie détectée & Oui ($\rho = 1.25$) \\
  Arcs filtrés (passabilité) & 3 arcs résidentiels exclus \\
  Raison du recalcul & Minuterie 5 min \\
  Chemin précédent & [N1043 $\to$ B5 $\to$ N8] \\
  Gain estimé & +2 min 40 s \\
  \bottomrule
\end{tabular}
\end{center}

\subsection{Schéma de présentation synthétique}

\begin{figure}[h]
\centering
\begin{tikzpicture}[
  layer/.style={rectangle, rounded corners=5pt, minimum width=10cm,
                minimum height=1.2cm, align=center, font=\small\bfseries},
  arr/.style={-{Latex[length=3mm]}, thick, gray}
]
  \node[layer, fill=urgence!15, draw=urgence] (l1) at (0,0)
    {Couche données : GPS, capteurs trafic, météo, signalements};
  \node[layer, fill=theorie!15, draw=theorie] (l2) at (0,-1.8)
    {Couche poids : mise à jour des multiplicateurs $\mu$, $\rho$, $\phi$};
  \node[layer, fill=attention!15, draw=attention] (l3) at (0,-3.6)
    {Couche algorithme : A* avec poids dynamiques + filtres passabilité};
  \node[layer, fill=pratique!15, draw=pratique] (l4) at (0,-5.4)
    {Couche décision : chemin optimal + logs d'audit + alertes};
  \node[layer, fill=gray!10, draw=gray] (l5) at (0,-7.2)
    {Couche conducteur : instructions GPS + override manuel};

  \draw[arr] (l1) -- (l2);
  \draw[arr] (l2) -- (l3);
  \draw[arr] (l3) -- (l4);
  \draw[arr] (l4) -- (l5);
\end{tikzpicture}
\caption{Architecture en couches du moteur de routage d'urgence.
  Chaque couche est indépendante et testable séparément.}
\end{figure}

\newpage

% ═══════════════════════════════════════════════════════════════════
%  SECTION 6 — SYNTHÈSE ET FEUILLE DE ROUTE
% ═══════════════════════════════════════════════════════════════════
\section{Synthèse et feuille de route}

\subsection{Récapitulatif des points clés}

\begin{center}
\begin{tabular}{>{\bfseries}p{3.5cm}p{9cm}}
  \toprule
  Concept & Signification opérationnelle \\
  \midrule
  Graphe pondéré & Modèle numérique du réseau routier avec temps de trajet \\
  Poids dynamique & Chaque arc a un coût qui évolue avec le trafic, la météo \\
  Vieillissement & Un chemin calculé il y a 5 min peut ne plus être optimal \\
  A* dynamique & Recalcul optimal toutes les 5 min (ou sur événement) \\
  Passabilité & Filtre selon le type de véhicule (ambulance $\neq$ camion) \\
  Invariants & Propriétés de sécurité garanties en toute circonstance \\
  Auditabilité & Chaque décision est loguée et explicable \\
  \bottomrule
\end{tabular}
\end{center}

\subsection{Extensions futures}

\begin{enumerate}
  \item \textbf{Apprentissage des multiplicateurs} — utiliser un modèle
        LSTM pour prédire $\mu(t+\Delta t)$ au lieu de l'interpolation
        linéaire ;
  \item \textbf{Routage multi-véhicules} — coordination de plusieurs
        ambulances pour éviter les conflits de ressources (arc surchargé
        par plusieurs véhicules d'urgence simultanément) ;
  \item \textbf{Intégration temps réel} — connexion directe aux APIs
        de trafic (HERE, TomTom, OpenStreetMap Overpass) pour des
        multiplicateurs à la minute ;
  \item \textbf{D* Lite} — remplacement de A* par D* Lite pour des
        recalculs encore plus rapides en ne retraitant que les arcs modifiés ;
  \item \textbf{Certification ISO 26262 / IEC 61508} — validation formelle
        des invariants de sécurité pour homologation dans un système
        critique embarqué.
\end{enumerate}

% ═══════════════════════════════════════════════════════════════════
%  CONCLUSION
% ═══════════════════════════════════════════════════════════════════
\section*{Conclusion}
\addcontentsline{toc}{section}{Conclusion}

Le moteur de routage d'urgence présenté dans ce document repose sur des
fondements algorithmiques solides (A* avec garantie d'optimalité) et une
modélisation réaliste des contraintes opérationnelles (trafic, météo,
passabilité, délais). La stratégie de \textbf{vieillissement des poids
toutes les 5 minutes} assure que le système reste aligné avec la réalité
du terrain, tandis que les \textbf{invariants métier} garantissent la
sécurité en toute circonstance.

\medskip
La conception est \textbf{transparente} : chaque décision est loguée,
chaque paramètre est explicable, et le conducteur conserve toujours la
primauté sur les suggestions du système. Ce n'est pas un oracle opaque
— c'est un \emph{assistant calculateur} au service des équipes d'urgence.

\medskip
Les données réelles de Yaoundé (2\,400 relevés de trafic sur 6 types de
voies, multiplicateurs allant jusqu'à $\times 3{,}33$) confirment la
pertinence de l'approche : les conditions de circulation urbaine sont
suffisamment variables pour justifier un recalcul fréquent, et suffisamment
structurées pour être modélisées efficacement par l'algorithme proposé.

\vspace{2em}
\begin{center}
{\color{urgence}\rule{0.5\linewidth}{1pt}}\\[0.5em]
{\itshape \og{}L'optimal d'aujourd'hui est le sous-optimal de demain.\\
Le recalcul n'est pas une faiblesse — c'est la définition de l'intelligence adaptative.\fg{}}\\[0.5em]
{\color{urgence}\rule{0.5\linewidth}{1pt}}
\end{center}

\end{document}